\scanpage{66}
\mspage{10}
% archive sheet number: 65
où $C:\mathrm{Cat}_{/S}\to\underline{\unsure{\mathrm{Parf}|\mathrm{Perf}}}\,S$,
\inserted{$\Phi$ et $\Psi$} sont les 3 foncteurs \inserted{``rétrogrades''} \unsure{\deleted{ment}ionnés|mentionnés} plus haut, et $\unsure{r|i}:\underline{\mathrm{Parf}}\,S\hookrightarrow\mathrm{Cat}_{/S}$ est l'inclusion.
Il faut donc vérifier que $\Phi_0,\Psi_0$ définis
par $(**)$, qui à priori \deleted{\unsure{prennent}} ont comme catégories but $\underline{\mathrm{Fibsc}}\,S$ et $\underline{\mathrm{Cofib}}\,\underline{\mathrm{sc}}\,S$, prennent \unsure{leurs|leur}
valeurs dans les sous-catégories avec indices $0$.

Je vais donner un énoncé d'\unsure{équivalence|équivalences}
unique, qui couvrira \inserted{(presque)} tous les \deleted{\unsure{8}} \deleted{\unsure{18}}16 cas à la
fois :

\uline{Corollaire} Considérons les diagrammes
\diagram{\begin{tikzcd}[ampersand replacement=\&,column sep=small,row sep=small]
\&\underline{\mathrm{Parf}}\,S\arrow[dd,"r"']\&\\
\underline{\mathrm{Fibsc}}_0\,S\arrow[dr,"p"']\&\&\underline{\mathrm{Cofibsc}}_0\,S\arrow[dl,dash,"q"]\\
\&\mathrm{Cat}_{/S}\&
\end{tikzcd}}
et soit $\deleted{\unsure{(}}\unsure{\mathcal F|F}\deleted{\unsure{,\Sigma)}}$ une \inserted{ss-}catégorie \inserted{pleine de} \deleted{\unsure{avec un} ``localisateur''
$\Sigma\subset\mathrm{Fl}(\mathcal F)$, et $f:\mathcal F\to\mathrm{Cat}_{/S}$ un foncteur \deleted{\unsure{pleinement fidèle}} compatible avec les localisateurs. Supposons qu'il
existe un $\mathrm{Cat}_{/S}$-foncteur $g$ soit de $\underline{\mathrm{Fibsc}}_0\,S$, soit
de $\underline{\mathrm{Cofibsc}}_0\,S$, \inserted{soit de $\underline{\mathrm{Parf}}\,S$,} dans $\mathcal F$, compatible avec les
localisateurs. Alors $f$ et $g$ induisent des
équivalences entre les catégories localisées.}
\inserted{$\mathrm{Cat}_{/S}$, telle que l'un des foncteurs $p,q,r$ se factorise par $\mathcal F$. Soit $W_{\mathcal F}$ induit sur $\mathcal F$ par $W_S$ de $\mathrm{Cat}_{/S}$. Alors
$\mathcal F\hookrightarrow\mathrm{Cat}_{/S}$ induit une équivalence sur les catégories}
\deleted{(Ce \deleted{\illegible{1}} \inserted{cor.} implique \inserted{déjà} que $p$ \deleted{\unsure{et $f$}} et $q$ induisent des équivalences,
comme on voit en prenant $\mathcal F=\mathrm{Cat}_{/S}$, $f=\text{foncteur identique}$.)}
\inserted{localisées, et il en est de $=$ de $p,q,r$.}

\deleted{\uline{\unsure{Dém.}} On va se borner au cas $\underline{\mathrm{Fibsc}}_0\,S$, et
il suffira de prouver que $f$ induit une équivalence.
Pour conclure qu'il en sera de même de $g$,
il suffit de prouver qu'il en est ainsi de $p$.
Mais ce sera le cas, en appliquant le cor. au cas
$\mathcal F=\underline{\mathrm{Fibsc}}_0\,S,\ \Sigma=W_S,\ f=p$,}
